% ============================================================
% SECȚIUNEA 1 — Descrierea arhitecturii
% Fișiere sursă de referință:
%   backend/AIInterviewCoach.Domain/
%   backend/AIInterviewCoach.Application/
%   backend/AIInterviewCoach.Infrastructure/
%   backend/AIInterviewCoach.API/
%   backend/AIInterviewCoach.API/Extensions/ServiceCollectionExtensions.cs
%   backend/AIInterviewCoach.API/Program.cs
% ============================================================

\section{Arhitectura sistemului}
\label{sec:arhitectura}

\subsection{Principii de proiectare}

Aplicația back-end a fost structurată conform paradigmei \emph{Clean Architecture}
(arhitectură curată), propusă de Robert C.\ Martin, care impune separarea
responsabilităților în straturi concentrice independente de tehnologie.
Dependențele de cod sunt permise exclusiv spre interior: straturile exterioare
pot referi straturile interioare, dar niciodată invers.
Avantajul practic este că logica de business poate fi testată izolat, fără
nicio dependență de baze de date, cadre web sau servicii externe.

\subsection{Stratificarea proiectului}

Soluția conține patru proiecte \texttt{.NET~10}, fiecare corespunzând unui
strat al arhitecturii:

\paragraph{AIInterviewCoach.Domain}
Cel mai interior strat; nu depinde de nicio altă bibliotecă din proiect.
Conține entitățile de domeniu (\texttt{User}, \texttt{Problem}, \texttt{Submission},
\texttt{Interview} etc.), enumerările (\texttt{UserRole}, \texttt{SubmissionStatus},
\texttt{InterviewSessionStatus}) și constantele de domeniu
(\texttt{SubmissionFeedbackStatuses}).
Aceste clase nu moștenesc din nicio clasă de infrastructură și nu referă
Entity Framework, ASP.NET sau orice alt cadru extern.
% Sursă: backend/AIInterviewCoach.Domain/Entities/ și Enums/

\paragraph{AIInterviewCoach.Application}
Stratul de aplicație orchestrează logica de utilizare.
Definește interfețele de servicii (\texttt{IAuthService}, \texttt{ISubmissionService},
\texttt{IProblemService} etc.) din directorul
\texttt{Application/Interfaces/Services/},
implementează serviciile concrete (\texttt{AuthService}, \texttt{SubmissionService},
\texttt{InterviewService}, \texttt{SubmissionFeedbackProcessor} etc.)
și declară obiectele de transfer de date (DTO-urile) din subdirectorul
\texttt{Application/DTOs/}.
Accesul la baza de date se face exclusiv prin interfața \texttt{IAppDbContext},
injectată prin constructor, nu prin clasa concretă \texttt{AppDbContext}.
% Sursă: backend/AIInterviewCoach.Application/

\paragraph{AIInterviewCoach.Infrastructure}
Stratul de infrastructură furnizează implementările concrete ale
interfețelor definite în Application.
Conține:
\begin{itemize}
  \item \texttt{AppDbContext} — contextul Entity Framework Core 10,
        configurat pentru PostgreSQL (provider implicit) sau SQLite
        (pentru teste de integrare), conform cheii de configurare
        \texttt{DatabaseProvider};
        % Sursă: ServiceCollectionExtensions.cs, AddApplicationDatabase()
  \item \texttt{JwtTokenService} — generează tokenuri JWT semnate HMAC-SHA256
        cu pretenții \texttt{sub}, \texttt{email}, \texttt{name} și \texttt{role};
        % Sursă: Infrastructure/Services/JwtTokenService.cs
  \item \texttt{PasswordHasherService} — hash-uiește parolele cu PBKDF2-HMAC-SHA256,
        10\,000 de iterații, salt aleatoriu de 16 octeți, stocând
        \texttt{base64(salt).base64(hash)};
        % Sursă: Infrastructure/Services/PasswordHasherService.cs
  \item \texttt{DotnetCodeExecutor} — execută cod C\#, Python și C++ în sandbox;
        % Sursă: Infrastructure/Services/DotnetCodeExecutor.cs
  \item \texttt{SubmissionFeedbackService} — apelează API-ul OpenAI;
        % Sursă: Infrastructure/Services/SubmissionFeedbackService.cs
  \item \texttt{SubmissionFeedbackBackgroundService} — coadă de fundal bazată
        pe \texttt{System.Threading.Channels}.
        % Sursă: Infrastructure/Services/SubmissionFeedbackBackgroundService.cs
\end{itemize}

\paragraph{AIInterviewCoach.API}
Stratul de prezentare expune API-ul REST prin ASP.NET Core Minimal
Controllers.
Conține controlere, middleware, politici de autorizare și rate limiting.
Controlarele sunt \emph{subțiri} în mod deliberat: ele validează
modelul de intrare (\texttt{ModelState.IsValid}), extrag identitatea
utilizatorului din claimuri JWT, deleghează operația serviciului de
aplicație corespunzător și returnează rezultatul.
Nu conțin logică de business.
% Sursă: API/Controllers/

\subsection{Fluxul unei cereri HTTP}

Parcursul complet al unei cereri tipice (exemplu: \texttt{POST /api/submissions})
este următorul:

\begin{enumerate}
  \item Cererea HTTP ajunge la pipeline-ul ASP.NET Core.
  \item \texttt{RequestObservabilityMiddleware} pornește un cronometru și
        înregistrează metrica la finalul cererii prin
        \texttt{IObservabilityService.RecordRequest()}.
        % Sursă: API/Middleware/RequestObservabilityMiddleware.cs
  \item \texttt{ExceptionHandlingMiddleware} înfășoară restul pipeline-ului
        într-un bloc \texttt{try/catch}; excepțiile negestionate sunt
        mapate la coduri HTTP: \texttt{KeyNotFoundException}→404,
        \texttt{UnauthorizedAccessException}→403,
        \texttt{InvalidOperationException}→400, altele→500.
        % Sursă: API/Middleware/ExceptionHandlingMiddleware.cs
  \item Middleware-ul de autentificare JWT validează cookie-ul \texttt{aic\_auth}
        sau antetul \texttt{Authorization: Bearer} și populează
        \texttt{HttpContext.User} cu claimurile tokenului.
        % Sursă: ServiceCollectionExtensions.cs, AddJwtAuthentication()
  \item Rate limiter-ul verifică politica \texttt{CandidateSubmissionFlow}
        (bucket de tokeni, 15 cereri/minut per utilizator autentificat).
        % Sursă: API/RateLimiting/RateLimitingPolicies.cs
  \item Middleware-ul de autorizare evaluează politica
        \texttt{CandidateWorkspaceAccess}, care cere rolul
        \texttt{Candidate} sau \texttt{Admin}.
        % Sursă: API/Authorization/AuthorizationPolicies.cs
  \item \texttt{SubmissionsController.CreateSubmission()} validează
        \texttt{ModelState}, extrage \texttt{candidateId} din claimul
        \texttt{ClaimTypes.NameIdentifier} și apelează
        \texttt{ISubmissionService.CreateSubmissionAsync()}.
        % Sursă: API/Controllers/SubmissionsController.cs
  \item \texttt{SubmissionService} interoghează baza de date prin
        \texttt{IAppDbContext}, execută codul candidatului prin
        \texttt{ICodeExecutor}, persistă submisia și pune ID-ul în coada
        de feedback AI prin \texttt{ISubmissionFeedbackQueue.QueueAsync()}.
        % Sursă: Application/Services/SubmissionService.cs
  \item Răspunsul DTO este serializat JSON și returnat cu HTTP~200.
\end{enumerate}

\subsection{De ce rămân controlarele subțiri}

Prin delegarea oricărei logici de business serviciilor din stratul Application,
proiectul obține mai multe beneficii:

\begin{itemize}
  \item \textbf{Testabilitate}: Serviciile pot fi testate cu baze de date
        InMemory (EF Core), fără a porni un server HTTP.
        Proiectul \texttt{AIInterviewCoach.Tests} conține 314 teste unitare
        și de integrare care nu folosesc nicio cerere HTTP reală.
  \item \textbf{Independență de transport}: Aceeași logică poate fi expusă
        ulterior prin gRPC sau mesagerie, fără a modifica serviciile.
  \item \textbf{Coeziune}: Un controler nu trebuie să cunoască regulile de
        acces la probleme (\texttt{ProblemAccessAuthorization}) sau
        algoritmul de hash al parolei; aceste detalii aparțin straturilor
        interioare.
\end{itemize}

\subsection{Înregistrarea serviciilor}

Toate serviciile sunt înregistrate în
\texttt{ServiceCollectionExtensions.AddApplicationServices()}:
singleton-uri pentru \texttt{IObservabilityService} și
\texttt{IProblemTemplateService} (fără stare mutabilă per cerere),
scoped pentru serviciile de aplicație care accesează baza de date,
și \texttt{IHostedService} pentru \texttt{SubmissionFeedbackBackgroundService}
(un singleton care implementează totodată \texttt{ISubmissionFeedbackQueue}).
% Sursă: API/Extensions/ServiceCollectionExtensions.cs, AddApplicationServices()
